Đồ án áp dụng mạng nơ-ron hồi tiếp (LSTM, GRU) và mạng tích chập một chiều (CNN) cho bài toán phát hiện và phân loại ransomware
trên bộ dữ liệu BitcoinHeist (2,9 triệu dòng, 7 lớp, lớp \code{white} chiếm 98,6\%). Vì dữ liệu ở dạng bảng, chúng tôi so sánh hai
cách biểu diễn chuỗi: (A) coi 11 đặc trưng của một dòng là một chuỗi 11 bước, và (B) dựng chuỗi lịch sử thật của mỗi địa chỉ
qua các ngày (tối đa 16 lần xuất hiện), đánh giá trên cách chia theo địa chỉ để tránh rò rỉ. Hơn 40 cấu hình được thử nghiệm có
hệ thống (số tầng, số units, hai chiều, cách gộp chuỗi, BatchNorm, Dropout, kích thước kernel, pooling, learning rate, batch size,
độ dài lịch sử), các cấu hình chính chạy với 3 seed. Với biểu diễn A, RNN/CNN đạt macro F1 0,51--0,54 --- ngang MLP và thấp hơn
XGBoost (0,575): xếp đặc trưng thành chuỗi không tạo thêm thông tin. Với biểu diễn B, macro F1 tăng lên 0,72--0,74; mô hình
nơ-ron tốt nhất là CNN kernel 5 (0,742), vượt XGBoost cùng thông tin một dòng (0,565) nhưng vẫn thấp hơn XGBoost có đặc trưng lịch sử
tổng hợp (0,767). Phân tích chi tiết cho thấy RNN/CNN vượt XGBoost trên các địa chỉ có lịch sử dài ($\ge$6 lần xuất hiện) nhưng
thua trên các địa chỉ mới --- nhóm chiếm 90\% dữ liệu. Biểu diễn B dễ overfit hơn và được lợi rõ rệt từ BatchNorm + Dropout,
trong khi với biểu diễn A điều chuẩn thêm chỉ gây underfit.
